\newpage

\phantomsection
\addcontentsline{toc}{chapter}{List of Publications}
\chapter*{Publications}

% %%%%%%%%%%%%%%%%%%%%%%%%%%%%%%%%%%%%%%%%%%%%%%%%%%%%%%%%%%%%%%%%%%%%%%%%%%
% Replace the statement below with full references (authors, title, venue,
% year, status) as manuscripts are submitted or accepted.
% %%%%%%%%%%%%%%%%%%%%%%%%%%%%%%%%%%%%%%%%%%%%%%%%%%%%%%%%%%%%%%%%%%%%%%%%%%

No publications have been submitted from this work at the time of thesis
submission. The following manuscripts, based on the results reported here, are in
preparation for submission to \textit{Physical Review~A} or \textit{Quantum
Science and Technology}.

\section*{Manuscripts in preparation}

\begin{enumerate}
    \item \textit{Matched classical controls for learned quantum error mitigation:
    classical kernel ridge regression matches or outperforms a fidelity quantum
    kernel at full training size.}
    Reports the ten-method benchmark and circuit-grouped evaluation protocol of
    Chapter~\ref{chap5:results}, the matched classical RBF kernel control, the
    repeated-split and initialisation-seed study
    (Section~\ref{sec:ch5_robustness}) and the depth-extrapolation split, on
    which the fidelity kernel falls below no mitigation while classical kernel
    ridge does not.

    \item \textit{Non-Markovian memory and circuit cost in learned error
    mitigation.}
    Reports the device-informed noise regime, the reset control that separates
    the defect's memory from its energy loss
    (Section~\ref{sec:ch5_markovianity}), and the circuit-execution cost of
    shot-estimated fidelity kernels (Section~\ref{sec:ch5_shotkernel}).
\end{enumerate}
